\begingroup
\renewcommand{\arraystretch}{1.45}
\fancyhead[L]{\small ZUR THEORIE DER NEUEN TRANSZENDENTEN. V.}

\Needspace{22\baselineskip}
\par\medskip\Needspace{5\baselineskip}\pdfbookmark[1]{Art. 16}{transscendentes-v-art-16}\centerline{[16.]}\nobreak\smallskip
Wenn innerhalb einer begrenzten Figur überall
\[
\frac{\mathrm{dd}V}{\mathrm{d}x^2}
+\frac{\mathrm{dd}V}{\mathrm{d}y^2}=0
\]
an der Grenze hingegen $V$ konstant $=A$ ist, so ist notwendig auch im ganzen Raum $V=A$


\Needspace{20\baselineskip}
Beweis. Es sei $\mathrm{d}M$ ein Element der Fläche und
\[
\begin{gathered}\int\Bigl\{
\Bigl(\frac{\mathrm{d}V}{\mathrm{d}x}\Bigr)^2
+\Bigl(\frac{\mathrm{d}V}{\mathrm{d}y}\Bigr)^2
\\[8pt]\quad{}+\Bigl(\frac{\mathrm{dd}V}{\mathrm{d}x^2}
+\frac{\mathrm{dd}V}{\mathrm{d}y^2}\Bigr)(V-A)
\Bigr\}\mathrm{d}M=\Omega\end{gathered}
\]


\Needspace{25\baselineskip}
\SourceLeaf{489}

Wäre $V$ nicht konstant, so wäre offenbar das Integral \emph{positiv}. Allein das Integral ist auch
\[
\int(V-A)\left(\frac{\mathrm{d}V}{\mathrm{d}x}\mathrm{d}y
-\frac{\mathrm{d}V}{\mathrm{d}y}\mathrm{d}x\right)
\]
durch den Umfang der Figur ausgedehnt, also $=0$. Da dies mit dem Vorigen
im Widerspruch steht, so ist die Voraussetzung unzulässig.

Ähnliches findet mit drei Unbestimmten statt.

Es lassen sich hieraus manche schöne Folgerungen ziehen.

Ein Punkt kann innerhalb eines hohlen Raumes nicht im stabilen Gleichgewicht sein, wenn nicht innerhalb des ganzen Raums gar keine Wirkung statt
findet, weder im Fall der Abstoßung, noch der Anziehung.


\Needspace{20\baselineskip}
\par\medskip\Needspace{5\baselineskip}\pdfbookmark[1]{Art. 17}{transscendentes-v-art-17}\centerline{[17.]}\nobreak\smallskip
Dass in der Peripherie einer Gleichgewichtsfigur keine \emph{negativen} Teile sein
können, beweist sich so. Gesetzt $AB$ wäre negativ, so wäre $V-A$ ($A$ Wert
von $V$ in der Peripherie) außen neben $AB$ positiv, neben den anderen Teilen negativ. Es wird daher einen endlichen Raum geben, wo dieses Positive gilt;
er sei innerhalb $ABCDE$; dann wäre er aber, da er an der Grenze $0$ ist, in
diesem Raum konstant, was ein Widerspruch ist.

Dieselbe Schlussart lässt sich auf drei Dimensionen anwenden.


\Needspace{18\baselineskip}
Bei drei Dimensionen findet, für die Fläche $s$, wo $V=\mathrm{Const.}$, noch das
schöne Theorem statt, dass $\int\frac{\mathrm{d}V}{\mathrm{d}\rho}\mathrm{d}s$,
wenn $\rho$ normal gegen $s$, $=4\pi M$, wo $M$
ganze Masse (vermutlich allemal die in einer solchen Fläche eingeschlossene
Masse, also ungerechnet die draußen liegende).

Vermutlich ist der Satz für \emph{jede} einhüllende Fläche gültig, ohne dass $V$
konstant zu sein braucht.


\Needspace{20\baselineskip}
Es seien zwei einander einschließende Flächen, in denen $V$ konstant ist,
nämlich jeweils $V=A$, $V=B$, $\mathrm{d}W$ ein Element des Raumes zwischen beiden,
$p$ die Anziehungskraft in jedem Punkt, dann ist
\[
\int pp\,\mathrm{d}W=(B-A)\, .\,4\pi M
\]
wenn $M$ die Masse im inneren Raum ist, wobei angenommen ist, dass im Raum
$W$ keine anziehende Masse ist.

\par\nopagebreak\begin{center}\rule{93.96bp}{0.35pt}\end{center}

\clearpage
\endgroup
